% Generated by Listas/make_statements.py from pratica-lecture-01.md, item 1. Do not edit.
\textbf{Source.} Tadelis (2013), \emph{Game Theory: An Introduction}, Exercise 3.2 ``Penalty Kicks'', p.~58 {[}PDF 77{]}. Reproduced verbatim. Answer in the source: none in the book (the preface, p.~xiii, says about half are solved on Princeton UP's website; not checked).



\begin{quote}

\textbf{Penalty Kicks:} Imagine a kicker and a goalie who confront each other in a penalty kick that will determine the outcome of a soccer game. The kicker can kick the ball left or right, while the goalie can choose to jump left or right. Because of the speed of the kick, the decisions need to be made simultaneously. If the goalie jumps in the same direction as the kick, then the goalie wins and the kicker loses. If the goalie jumps in the opposite direction of the kick, then the kicker wins and the goalie loses. Model this as a normal-form game and write down the matrix that represents the game you modeled.

\end{quote}

